\documentclass{article}
\usepackage{amsmath,amssymb,hyperref}
\title{A local veto requires both distinguishing evidence and a gate that can act on it}
\begin{document}
\maketitle

A lock can be asserted without excluding a particular command. That can happen even when the observer, topic, and lock value are unchanged. The missing qualification is which command the gate can veto. The selected ROS command multiplexer makes this distinction exact, and also limits what a single ordered priority table can express across modes.

\paragraph{What the receiver actually does.}
At the \href{https://github.com/ros-teleop/twist_mux/tree/a8225e2d9bf0171aac447c246376f0981defd17b}{pinned \texttt{twist\_mux} source}, a lock counts as asserted when its last Boolean is true or its positive timeout has expired. Expiry is strictly later than the timeout since the local callback stamp. A zero timeout disables expiry. The common threshold is the maximum priority of the asserted locks, or zero if none. The selected unstamped velocity path masks a command with priority strictly below the threshold; equality is not masked. It then chooses the highest-priority unmasked, unexpired source, breaking ties by encounter order in the selected loop. The callback updates the arriving source's stamp before evaluating that choice. A lock callback records its new value and time; it does not itself publish a stop from this receiver. These are facts about command selection, not a guarantee that an actuator is stopped or no longer occupies a shared space.

The replay compiles those exact selected functions with stubs for messages, time, and publishing. An active lock of priority two affects an otherwise selectable command of priority one but does not affect a command whose priority is two. The source priority and scope belong in any transported veto certificate. Receipt, activation of a lock, effective veto of a named action, and physical stopping are different observations. No communications security fault or live device test is involved.

A separate replay stops publication from a higher-priority incumbent, then submits lower-priority requests from an otherwise eligible replacement. The receiver keeps selecting the incumbent while its last command is current. This suppresses the replacement's callbacks without republishing the incumbent's packet. If the incumbent stops publishing at time zero and its timeout is \(T>0\), the first eligible replacement publication is its first admissible callback strictly after \(T\), not the logical controller-retirement time. With zero timeout the old eligibility can persist indefinitely in the fixed fixture. A scalar lock cannot remove just the higher-priority incumbent while keeping the lower-priority replacement. These facts concern the receiver's cached eligibility; nothing in the replay determines whether a downstream drive holds its earlier command. When reasoning about the \href{https://github.com/eternal-november/eos-research-programmes/pull/460}{strict lifecycle switch}, this is a possible additional obligation of a separate command layer. The two upstream projects are not assumed to be wired together.

\paragraph{An exact architectural restriction.}
Fix static priorities \(p_j>0\) for a finite set of sources. For clarity evaluate each source in isolation, with its command current and all competing commands inactive. For threshold \(t\), the allowed set is \(\{j:p_j\ge t\}\). For any pair of thresholds, the allowed sets are nested: raising the threshold can only remove elements. Conversely, any chain of required allowed sets under inclusion can be realized by some fixed priority ordering and choice of thresholds if there are enough priority levels, including enough room to block everything when required. Order elements by the layer where they first enter the increasing chain, putting early entrants at high priority; subsequent thresholds can then select exactly those sets. A priority table need not use distinct priorities for elements that always enter together.

In particular a single threshold cannot realize both the exact permission \(\{A\}\) in one mode and \(\{B\}\) in a different mode with the same priority ordering, even if every observer always knows the mode. This obstruction comes from available control, not observability. A different actuator-side interface with mode-dependent command sets or independently scoped gates is a real change of the control architecture. It is not equivalent to transmitting a more detailed lock report to the original single threshold. This characterization does not claim simultaneous useful service when several commands compete for the one output.

Several independently set thresholds, with a command allowed only when it passes them all, give intersections of these ordered permission relations. This is the established \href{https://www.sciencedirect.com/science/article/pii/S0167506008708896}{Ferrers dimension} of the mode--command incidence relation, with the physical interpretation that each dimension is an independently enforceable threshold. There are two useful exact extremes. If each mode allows just one of \(n\ge2\) commands, two axes suffice by using opposite orderings: a lower and an upper rank bound select the desired command. One axis is impossible by nesting. If instead each mode only forbids its own command and allows all others, \(n\) axes are necessary. For one axis to be responsible for excluding command \(j\) in its mode while retaining every other command, \(j\) must be strictly below every other command on that axis. One axis cannot do this for two different commands. Giving each command such a dedicated lowest-rank axis attains the bound, assuming enough usable ranks. The result is about this modeled gate family, not an impossibility for an addressed deny mask or a rank table that can lawfully change with the mode. A controllable, timely hardware path is still required for each axis.

\paragraph{Fair supervisory-control comparison.}
For arbitrary qualified gates, compare with \href{https://xiangyin.sjtu.edu.cn/Paper/15Automatica.pdf}{Yin and Lafortune, Section 2.2, Definition 2}. At a single cut, the declared finite plant histories each have a candidate command marked required allowed or required forbidden. Every history supplies local observation keys and the actually effective, timely subset of actors with a veto on this action. Those are plant inputs, not inferred from an actor's knowledge. Allow an extra forced veto bit for a physical default even if no person can issue a timely choice.

No forced veto may occur in a required allowed history. Protect each local key used by an effective actor in such a history: that actor must not veto on that key. Every required forbidden history must then have either a forced veto or an effective actor with an unprotected key. Necessity follows from the allowed histories; sufficiency follows by letting actors veto on exactly unprotected keys. This is a cut-local conjunctive supervisor condition with concrete control and timeout premises, not a new language-synthesis theorem. An actor seeing the right condition but unable to affect the gate in time contributes no veto. Conversely, permitting an action after an unavailable observer's timeout is impossible in this interface if the forced default already excludes it. A later authenticated observation cannot change the completed cut.

The exact predicate ties full local table search in 4,096 cases. Direct ordering enumeration matches the nested-family test in 2,696 cases, verifies the two multi-axis extremes through six commands, and a direct threshold and timing calculation agrees with 18,432 native selected callback cases and 720 retired-producer replay cases. Lean checks nesting and the two-mode obstruction. No full ROS, physical clearance, independent non-radio source substitution or research-contract closure is claimed.
\end{document}
